% TOOL:    papers/_content/ldpe-blown-films-lamination.ieee.tex
% FAMILY:  papers
% VERSION: 1.0.0
% DATE:    2026-11-13
% CHAT:    LDPE blown films for lamination white paper
% CHANGES: 1.0.0 - IEEE-style two-column A4 source for Technical Paper No. 05; build: pdflatex (run twice)
% STATUS:  working
\documentclass[10pt,twocolumn,a4paper]{article}
% IEEE conference-style layout (A4), emulated with article class
\usepackage[utf8]{inputenc}
\usepackage[T1]{fontenc}
\usepackage{mathptmx}
\usepackage{textcomp}
\usepackage{newunicodechar}
\usepackage[a4paper,top=19mm,bottom=43mm,left=14.32mm,right=14.32mm,columnsep=4.22mm]{geometry}
\usepackage{amsmath}
\usepackage{array,tabularx,booktabs}
\usepackage{titlesec}
\usepackage[font=footnotesize,labelsep=newline,justification=centering,singlelinecheck=false]{caption}
\usepackage{cite}
\usepackage{xurl}
\usepackage[hidelinks]{hyperref}
\urlstyle{same}

\newunicodechar{µ}{\textmu}
\newunicodechar{°}{\textdegree}
\newunicodechar{≥}{\ensuremath{\geq}}
\newunicodechar{≤}{\ensuremath{\leq}}
\newunicodechar{≈}{\ensuremath{\approx}}
\newunicodechar{±}{\ensuremath{\pm}}
\newunicodechar{×}{\ensuremath{\times}}
\newunicodechar{₂}{\ensuremath{_2}}
\newunicodechar{₃}{\ensuremath{_3}}
\newunicodechar{³}{\ensuremath{^3}}
\newunicodechar{²}{\ensuremath{^2}}

% IEEE-style headings
\renewcommand{\thesection}{\Roman{section}}
\renewcommand{\thesubsection}{\Alph{subsection}}
\titleformat{\section}{\centering\normalfont\scshape}{\thesection.}{0.5em}{}
\titlespacing*{\section}{0pt}{8pt plus 2pt}{4pt}
\titleformat{\subsection}{\normalfont\itshape}{\thesubsection.}{0.5em}{}
\titlespacing*{\subsection}{0pt}{6pt plus 2pt}{3pt}

% IEEE-style table captions: "TABLE I" over small-caps title
\renewcommand{\thetable}{\Roman{table}}
\captionsetup[table]{labelfont={normalfont},textfont={sc},name=TABLE,skip=3pt}

\setlength{\parindent}{1em}
\setlength{\parskip}{0pt}
\renewcommand{\arraystretch}{1.15}
\newcolumntype{L}{>{\raggedright\arraybackslash}X}
\pagestyle{empty}

\begin{document}

\twocolumn[{%
\begin{center}
{\fontsize{24}{28}\selectfont LDPE-Based Blown Films for Lamination:\\[2pt] Historical Emergence, Process Windows, Additives and Current Trends\par}
\vspace{14pt}
{\large K Vijay Anand\par}
\vspace{4pt}
{\small Technical Paper No.~05 \quad\textbullet\quad \url{https://haivijayanand.github.io/papers/}\par}
\vspace{12pt}
\end{center}
}]

\noindent\textbf{\textit{Abstract}---LDPE-based blown film remains the workhorse sealant web of laminated flexible packaging, because long-chain branching gives it the melt strength, bubble stability and clarity that lamination converters need. Over nine decades the sealant has moved from monolayer LDPE at 80–100 µm to 5–9 layer co-extrusions of LDPE, LLDPE and metallocene PE at 30–60 µm, with sealing starting 20–30 °C lower and dart impact several times higher. This paper traces that history, sets out the process window (blow-up ratio, draw-down, frost line and the temperature at each station of the line), summarises the property requirements of lamination-grade film, and gives a separate review of additives: metallocene resins, fillers, impact modifiers, slip, antiblock, processing aids and stabilisers. It closes with current trends: all-PE recyclable laminates, recycled content and PFAS-free processing aids. Numerical ranges are typical industry values for guidance; resin supplier data sheets and line trials govern any specific set-up.}

\vspace{6pt}
\noindent\textbf{\textit{Index Terms}---Blown film extrusion, LDPE, lamination, blow-up ratio, metallocene polyethylene, polymer additives, flexible packaging, recyclability.}

\section{Introduction}

In a laminate, the PE film is the inside web: it seals the pack, carries the product contact and provides most of the toughness, while the outer web (BOPET, BOPP, BOPA, paper or aluminium foil) carries print, stiffness and barrier. Typical structures are PET 12 µm // PE 40–100 µm for snacks, spices and liquids, PET // Al foil // PE for high barrier, and BOPP // PE or MDO-PE // PE for all-PE recyclable packs.

The sealant film must therefore do four jobs at once:
\begin{itemize}\setlength\itemsep{1pt}
\item \textit{Bond}: accept adhesive or extrudate with a stable surface energy (≥ 40 dyn/cm at lamination) and no migrating additives at the interface.
\item \textit{Seal}: start sealing at low temperature, with high hot tack and seal strength through contamination.
\item \textit{Survive}: resist puncture, flex-cracking and drop impact in the filled pack.
\item \textit{Run}: have controlled COF, flatness, gauge uniformity and low gels on the laminator and the form-fill-seal (FFS) machine.
\end{itemize}
LDPE alone does the first and last well; modern film design adds LLDPE and metallocene PE to do the middle two.

\section{Historical Emergence}

LDPE film went from a wartime curiosity to the default sealant in about 30 years \cite{fawcett,cantor}; every later step has been about doing the same job with less material and lower sealing energy. Table~\ref{tab:history} summarises the milestones.

\begin{table*}[!t]
\caption{Milestones in LDPE-Based Lamination Film}
\label{tab:history}
\footnotesize
\begin{tabularx}{\textwidth}{@{}p{24mm}LL@{}}
\toprule
\textbf{Period} & \textbf{Development} & \textbf{What it changed for lamination film}\\
\midrule
1933 & ICI (Fawcett and Gibson) obtain polyethylene by high-pressure polymerisation of ethylene & A tough, inert, heat-sealable polymer becomes possible\\
1939–1945 & Commercial LDPE at ICI; used mainly for radar cable insulation & Capacity built; post-war surplus looks for markets\\
Late 1940s–1950s & Tubular blown film process adopted for LDPE & Thin, wide, low-cost film; bags and liners\\
1950s–1960s & Ziegler and Phillips catalysts give HDPE; extrusion coating and early adhesive lamination of cellophane, paper and foil & PE becomes the heat-seal layer of laminates\\
1970s & BOPP and BOPET films grow; polyurethane adhesive lamination spreads; gas-phase LLDPE (Unipol, 1977) & Two-web PET//PE and BOPP//PE laminates take over; LLDPE adds toughness\\
1980s & 3-layer co-extrusion; LDPE/LLDPE blends; C6 and C8 comonomers & Sealant, core and treat-side layers tuned separately\\
1990s & Single-site metallocene PE and plastomers (Exxon Exact 1991; Dow Insite/Affinity/Elite 1993–94); solventless adhesives & Low seal initiation, high hot tack, downgauging\\
2000s & 5–7 layer lines, internal bubble cooling, automatic gauge control; EVOH/PA barrier co-ex & Faster lines, ±3–5\% gauge, barrier sealants\\
2010s & 9–11 layer lines; machine-direction oriented (MDO) PE & Stiff, printable PE replaces PET as outer web\\
2020s & Design-for-recycling guidelines (CEFLEX, APR); EPR and recycled-content rules, including India's Plastic Waste Management Amendment Rules 2022; PFAS-free processing aids & All-PE laminates, PCR in core layers, additive reformulation\\
\bottomrule
\end{tabularx}
\end{table*}

The common thread is that LDPE never left the film. It moved from being the whole film to being the 10–40\% melt-strength component that keeps the bubble stable while LLDPE and metallocene PE carry the mechanical and sealing duty.

\section{The Blown Film Process and Blow-Up Ratio}

For lamination-grade LDPE-rich film, a blow-up ratio (BUR) of 2.0–3.0 with the frost line at 2–4 die diameters gives the best balance of MD/TD properties, flatness and bubble stability \cite{cantor}.

Resin is melted in one or more single-screw extruders, filtered, and fed through an adapter into an annular spiral-mandrel die. The emerging tube is inflated with air, cooled by an air ring (and on modern lines by internal bubble cooling, IBC), and solidifies at the frost line. It is then flattened in a collapsing frame, pulled by the primary nip, corona-treated on the bonding side, and wound. Orientation is set between the die lip and the frost line by two independent stretches: radial (blow-up) and axial (draw-down).

\subsection{Key Geometric Ratios}

The blow-up ratio, layflat width and draw-down ratio (DDR) are
\begin{equation}
\mathrm{BUR} = \frac{D_{\text{bubble}}}{D_{\text{die}}}
\end{equation}
\begin{equation}
W_{\text{layflat}} = \frac{\pi D_{\text{die}}\,\mathrm{BUR}}{2} \approx 1.571\,D_{\text{die}}\,\mathrm{BUR}
\end{equation}
\begin{equation}
\mathrm{DDR} = \frac{h_{\text{die gap}}}{t_{\text{film}} \times \mathrm{BUR}}
\end{equation}
where $D$ is diameter, $h$ the die gap and $t$ the film thickness. For example, a 250 mm die at BUR 2.5 gives a 625 mm bubble and a 982 mm layflat. With a 1.0 mm die gap and 50 µm film, DDR = 1000/(50 × 2.5) = 8.

\subsection{Typical Process Windows}

Table~\ref{tab:window} lists typical windows for monolayer LDPE and for co-extruded sealants.

\begin{table}[!t]
\caption{Typical Blown Film Process Windows}
\label{tab:window}
\footnotesize
\begin{tabularx}{\columnwidth}{@{}p{19mm}p{11mm}p{14mm}L@{}}
\toprule
\textbf{Parameter} & \textbf{LDPE mono} & \textbf{LDPE/ LLDPE/mPE co-ex} & \textbf{Effect of increasing it}\\
\midrule
Blow-up ratio & 1.8–3.0 & 2.0–3.5 & More TD orientation, higher TD tear and dart, lower MD tear; above ≈3.5 bubble wanders\\
Die gap (mm) & 0.6–1.2 & 1.2–2.5 & Less shear and melt fracture; higher DDR and MD orientation\\
Draw-down ratio & 5–20 & 10–40 & More MD orientation, MD splittiness\\
Frost line (× die dia.) & 2–4 & 2–5 & Clearer, better-balanced film; lower output\\
Output (kg/h per cm die circ.) & 0.8–1.5 & 1.2–2.5 with IBC & Needs more cooling; raises gauge variation\\
Thickness for lamination (µm) & 40–120 & 30–100 & ---\\
\bottomrule
\end{tabularx}
\end{table}

Guidance for lamination webs:
\begin{itemize}\setlength\itemsep{1pt}
\item \textit{Balance over strength.} A BUR near 2.5 gives near-balanced MD/TD tear and shrink, which keeps the laminate flat and limits tunnelling and curl after curing.
\item \textit{Low BUR} (< 1.8) gives MD-oriented, splitty film with high MD shrink; it can tunnel on the laminator.
\item \textit{LDPE content sets the ceiling.} 10–30\% LDPE in an LLDPE/mPE blend is usually enough to hold a stable bubble at BUR 2.5–3.0.
\item \textit{Layflat is fixed by laminate width,} so BUR and die size are chosen together; a die change is often cheaper than running at an unstable BUR.
\end{itemize}

\section{Temperature Profile at Each Station}

LDPE lamination film is typically run with a rising barrel profile from about 160 °C to 200 °C and a melt temperature of 190–215 °C; LLDPE- and mPE-rich layers need 10–25 °C more. Table~\ref{tab:temp} gives typical set points station by station.

\begin{table*}[!t]
\caption{Typical Temperature Profile Along a Blown Film Line}
\label{tab:temp}
\footnotesize
\begin{tabularx}{\textwidth}{@{}p{50mm}p{30mm}p{32mm}L@{}}
\toprule
\textbf{Station} & \textbf{LDPE (MFI 0.5–2), °C} & \textbf{LLDPE / mLLDPE-rich, °C} & \textbf{Purpose and notes}\\
\midrule
Feed throat (water-cooled) & 30–50 & 30–50 & Stops bridging and premature melting in the hopper\\
Barrel zone 1 (feed) & 150–170 & 170–190 & Low enough to convey solids; too hot causes surging\\
Barrel zone 2 (transition) & 170–185 & 190–205 & Main melting zone\\
Barrel zone 3 (metering) & 180–195 & 200–215 & Homogenise melt, build pressure\\
Barrel zones 4–5 & 185–200 & 205–220 & Final melt conditioning\\
Screen changer / filter & 190–205 & 205–220 & 80–150 mesh pack holds back gels and contamination\\
Adapter and melt pipe & 190–205 & 205–225 & Equal to or slightly above barrel exit\\
Die body and spiral zones & 195–210 & 205–225 & Uniform around circumference; ±2 °C or better\\
Die lip & 195–215 & 210–230 & Hotter lips reduce sharkskin and die lines\\
Melt temperature (measured) & 190–215 & 205–235 & Higher gives clarity and fewer gels, but a weaker bubble\\
Air ring cooling air & 10–20 & 10–20 & Chilled air raises output; frost line moves down\\
IBC air (inlet) & 15–25 & 15–25 & Internal cooling adds 20–50\% output\\
Frost line (film) & ≈ 95–110 & ≈ 105–120 & Crystallisation; orientation frozen here\\
Collapsing frame / primary nip (film) & 40–60 & 40–60 & Too hot causes wrinkles, blocking, inside-sticking\\
Corona treater (film) & 30–50 & 30–50 & Treat to 42–48 dyn/cm for adhesive lamination\\
Winder (film) & < 35–40 & < 35–40 & Hot rolls block and lose dyne faster\\
\bottomrule
\end{tabularx}
\end{table*}

Practical points for lamination quality:
\begin{itemize}\setlength\itemsep{1pt}
\item \textit{Melt temperature drives clarity and gels.} Too cold leaves unmelt gels and haze; too hot oxidises the melt, adds carbonised gels and odour, and consumes antioxidant.
\item \textit{Treat-side temperature matters.} Corona treatment is most stable on film below 40 °C; dyne decays faster on warm, slip-loaded film.
\item \textit{Seasonal drift.} In a hot shop (35–40 °C ambient), chilled air holds the frost line; set points that work in winter can give a higher, wandering frost line in summer.
\item \textit{Co-extrusion.} Adjacent layers should differ by no more than about 20–30 °C at the die to avoid interfacial instability (zig-zag and wave patterns).
\end{itemize}
These figures are starting points; resin data sheets and a gel-count, dart and dyne trial set the final profile.

\section{Lamination-Grade Film Properties and Their Trend}

The biggest shift since the 1980s is in sealing and toughness: seal initiation has dropped from about 110 °C to 75–95 °C and dart impact per micron has risen three- to five-fold, which allowed sealant gauge to fall by roughly half.

\subsection{Typical Specification}

Table~\ref{tab:spec} lists a typical specification for a lamination sealant web.

\begin{table*}[!t]
\caption{Typical Specification for a Lamination Sealant Web}
\label{tab:spec}
\footnotesize
\begin{tabularx}{\textwidth}{@{}p{40mm}p{40mm}p{55mm}L@{}}
\toprule
\textbf{Property} & \textbf{Test method} & \textbf{Typical target} & \textbf{Why it matters in lamination}\\
\midrule
Thickness tolerance & ASTM D6988 / ISO 4593 & ±5\% (±3\% with auto-profile) & Gauge bands cause tunnelling, wrinkles, telescoping\\
Surface energy, treated side & ASTM D2578 & ≥ 40 dyn/cm at lamination (treat to 42–48) & Adhesive wetting and bond strength\\
COF, film/film (sealant side) & ASTM D1894 & 0.15–0.30 (auto-pack); 0.3–0.5 (manual) & Machinability on FFS; must hold after cure\\
Seal initiation temperature & ASTM F88 & 90–110 °C (LDPE); 75–95 °C (mPE) & Line speed, energy, sealing through PET\\
Hot tack & ASTM F1921 & ≥ 3–5 N/25 mm at seal temperature & Vertical FFS with heavy fills\\
Dart impact (F50) & ASTM D1709 A & 100–150 g (LDPE 50 µm); 300 to > 800 g (mPE blends) & Drop and puncture of filled packs\\
Haze & ASTM D1003 & < 8\% general; < 5\% clear packs & Reverse-printed clear windows\\
Gels (> 200 µm) & Optical gel counter & Low; specified per m² & Gel bumps show through PET and foil\\
Bond strength in laminate & ASTM F904 & > 2–3 N/15 mm (retort, liquids higher) & Delamination of pouches\\
Odour, migration & EU 10/2011; FDA 21 CFR 177.1520; IS 10146 & Pass & Food contact\\
\bottomrule
\end{tabularx}
\end{table*}

\subsection{Trend in Sealant Performance}

Table~\ref{tab:trend} shows how typical sealant performance has moved by era. Values are representative ranges compiled from typical resin data and converter practice, not a single data set. Other properties trended alongside: gel counts fell with better filtration and cleaner resins, gauge variation fell with automatic die-bolt or air-ring profile control, and treat-side dyne retention improved as slip moved to non-migratory systems (Section~\ref{sec:add}).

\begin{table}[!t]
\caption{Trend in Typical Sealant Performance}
\label{tab:trend}
\footnotesize
\begin{tabularx}{\columnwidth}{@{}p{11mm}Lp{10mm}p{10mm}p{13mm}@{}}
\toprule
\textbf{Era} & \textbf{Typical sealant} & \textbf{Gauge (µm)} & \textbf{SIT (°C)} & \textbf{Dart, 50 µm eq. (g)}\\
\midrule
1970s & Monolayer LDPE & 80–120 & 110–115 & 100–150\\
1980s–90s & LDPE/LLDPE (C4) blend, 3-layer & 60–90 & 105–110 & 150–300\\
2000s & 3–5 layer, C6/C8 LLDPE + mLLDPE skin & 50–70 & 95–100 & 300–600\\
2010s–20s & 5–9 layer, mLLDPE + plastomer skin, 10–30\% LDPE & 30–60 & 75–95 & 600 to > 1000\\
\bottomrule
\end{tabularx}
\end{table}

\section{Additives and Property Enhancements}
\label{sec:add}

In a lamination web every additive has two faces: what it does for the film, and what it does to the bond. The rule of thumb is to put functional additives only in the layer that needs them, and keep the treat side clean.

\subsection{Metallocene PE (mLLDPE) and Plastomers}

Single-site catalysts give PE with a narrow molecular-weight distribution ($M_w/M_n \approx$ 2–3) and comonomer spread evenly across all chains. This removes the low-molecular-weight, comonomer-rich fraction that makes Ziegler LLDPE sticky and hazy, and the high-density fraction that raises its sealing temperature.
\begin{itemize}\setlength\itemsep{1pt}
\item \textit{Grades:} mLLDPE (0.912–0.920 g/cm³) for toughness; plastomers/POP (0.875–0.905 g/cm³) for very low seal initiation and hot tack.
\item \textit{Typical use:} 20–60\% mLLDPE in core and sealant layers; 10–40\% plastomer in a 15–25\% sealant skin.
\item \textit{Gains:} dart impact two to five times Ziegler LLDPE at the same gauge; puncture and flex-crack resistance; seal initiation down by 10–30 °C; wider hot-tack window; seal-through-contamination; 20–40\% downgauging.
\item \textit{Costs:} low melt strength and high shear viscosity, giving higher motor load, melt fracture and a weaker bubble. The remedy is 10–30\% LDPE for melt strength plus a polymer processing aid (PPA). Plastomer-rich skins can raise COF and blocking, needing more slip and antiblock.
\end{itemize}

\subsection{LDPE as a Blend Partner}

LDPE is now often the additive rather than the base. 5–30\% LDPE in LLDPE or mLLDPE raises extensional viscosity and bubble stability, cuts haze through smoother surfaces, and lets the line run higher BUR and output. Above about 30–40\% LDPE, dart and tear begin to fall sharply.

\subsection{Fillers and Pigments}

Table~\ref{tab:fill} summarises common fillers. Fillers belong in the core layer of a co-extrusion, with clean, unfilled skins on both sides so that the treat and seal surfaces stay smooth.

\begin{table}[!t]
\caption{Fillers and Pigments in Lamination Film}
\label{tab:fill}
\footnotesize
\begin{tabularx}{\columnwidth}{@{}p{16mm}p{11mm}LL@{}}
\toprule
\textbf{Filler} & \textbf{Loading} & \textbf{Benefit} & \textbf{Penalty / caution}\\
\midrule
CaCO₃ MB (coated, 1–3 µm) & 5–20\% of layer & Cost, stiffness, faster cooling, dead-fold & Dart and tear drop; moisture causes voids and lacing (dry it); rougher surface\\
TiO₂ white MB (50–70\%) & 3–10\% MB & Opacity for white sealant, light-sensitive products & Die-lip build-up; abrasive; higher density\\
Talc & 1–5\% & Antiblock, stiffness, some nucleation & Haze; absorbs some additives\\
Carbon black MB & 2–6\% MB & UV and light barrier; black/white co-ex & Food-contact grade required\\
Desiccant (CaO) MB & 1–3\% & Removes moisture from recycled or CaCO₃-rich feed & Fixes moisture only, not contamination\\
\bottomrule
\end{tabularx}
\end{table}

\subsection{Impact Resistance and Toughness Enhancers}

\begin{itemize}\setlength\itemsep{1pt}
\item \textit{Comonomer choice:} C6 (hexene) and C8 (octene) LLDPE are markedly tougher than C4 (butene) at the same density.
\item \textit{mLLDPE and plastomers:} the most effective dart and puncture enhancers (Section~\ref{sec:add}-A).
\item \textit{EVA (4–18\% VA):} low-temperature toughness, flexibility and lower seal initiation; adds odour and blocking above about 9\% VA.
\item \textit{Ionomers:} best seal-through-contamination and hot tack; used as thin sealant skins in meat, cheese and liquid packs.
\item \textit{Layer design:} a stiff HDPE or MDPE core between tough mLLDPE layers gives both stiffness and impact; asymmetry must be balanced to avoid curl.
\end{itemize}

\subsection{Surface and Process Additives}

Table~\ref{tab:surf} lists the surface and process additives and their lamination cautions. Slip is the most frequent cause of bond problems: it migrates into the adhesive and to the treat side, lowering bond strength and dyne, and COF rises after cure as slip moves into the PET side.

\begin{table*}[!t]
\caption{Surface and Process Additives}
\label{tab:surf}
\footnotesize
\begin{tabularx}{\textwidth}{@{}p{45mm}p{30mm}p{45mm}L@{}}
\toprule
\textbf{Additive} & \textbf{Typical loading} & \textbf{Function} & \textbf{Lamination caution}\\
\midrule
Slip: erucamide (slow bloom), oleamide (fast bloom) & 300–1500 ppm in sealant skin & Lowers COF for FFS & Migrates to adhesive and treat side; use minimum, keep out of treat layer, or use non-migratory slip (siloxane MB, high-MW amides)\\
Antiblock: natural/synthetic silica, talc, CaCO₃ & 1000–4000 ppm & Prevents roll blocking; supports slip & Haze and roller abrasion; agglomerates read as gels\\
PPA: fluoropolymer, now PFAS-free alternatives & 300–1000 ppm active & Removes melt fracture with mLLDPE; cuts die build-up and gels & Needs conditioning; interacts with antiblock and TiO₂; for PFAS-free grades prefer silicone-free chemistries on laminated film and confirm dosing on trial\\
Antioxidants (hindered phenol + phosphite) & 500–2000 ppm & Protect melt at 200–230 °C; fewer carbonised gels & Over-oxidised film can yellow and give odour\\
Antistatic (GMS, ethoxylated amines) & 0.1–0.5\% & Less dust pick-up and static at winder & Migrates like slip; generally avoided in lamination webs\\
Antifog (glycerol/sorbitan esters) & 0.5–2\% & Clear view on chilled produce & Product side only; not in bonded layers\\
UV stabiliser (HALS) & 0.1–1\% & Outdoor and agricultural film life & Rarely needed in indoor laminates\\
Nucleating agents / clarifiers & 500–2000 ppm & Stiffness, lower haze, faster crystallisation & Shifts frost line; check seal window\\
Compatibilisers / odour absorbers & 1–5\% & Allow PCR in core layers & Verify gels, odour and food-contact status\\
\bottomrule
\end{tabularx}
\end{table*}

\subsection{Other Property Enhancements by Structure}

\begin{itemize}\setlength\itemsep{1pt}
\item \textit{Barrier:} an EVOH or PA core with maleic-anhydride tie layers gives oxygen barrier inside a PE structure; EVOH below about 5\% of the structure keeps it recyclable in PE streams \cite{ceflex}.
\item \textit{Stiffness and print face:} HDPE skins or MDO-PE (draw ratio 4–7) give PET-like stiffness and dimensional stability for all-PE laminates.
\item \textit{Heat resistance:} an HDPE or MDPE outer layer widens the sealing window and resists burn-through on fast jaws.
\item \textit{Downgauging:} mLLDPE toughness with an HDPE core lets 30–40 µm webs replace 60–80 µm LDPE/LLDPE films.
\end{itemize}

\section{Current Trends}

The PE sealant is becoming the whole laminate, built for recycling, with recycled content and without fluorinated processing aids.

\textit{All-PE laminates.} MDO-PE or BOPE replaces PET or BOPP as the printed outer web, laminated to a PE sealant. A published ExxonMobil and W\&H pouch uses a 25 µm MDO-PE // 120 µm PE laminate in place of PET12 // PE120 \cite{wh}. For the blown film line this means tighter gauge, flatness and MDO-ready resin in the outer web, while the LDPE-rich sealant stays largely unchanged.

\textit{Design-for-recycling limits.} CEFLEX D4ACE guidance (September 2025) focuses on mono-PE, mono-PP and mixed polyolefin flexibles and sets limits such as adhesives and coatings at ≤ 5\% by weight \cite{ceflex}. APR asks for testing when a PE film carries EVOH or PA \cite{pe2023}. Barrier therefore moves to thin EVOH cores, AlOx/SiOx coatings or metallised MDO-PE.

\textit{PFAS-free processing aids.} Tightening PFAS rules are pushing converters off fluoropolymer PPAs. Clariant launched silicone- and siloxane-free PFAS-free PPA grades in 2025, stating no effect on dyne, sealing or COF \cite{clariant}; Ampacet offers a PFAS-free PPA masterbatch at 1–3\% that it states does not interfere with printing or lamination \cite{ampacet}.

\textit{Recycled content.} PCR and PIR PE go into core layers, with virgin, clean skins on the treat and seal sides; this needs better melt filtration, desiccant masterbatch and odour control. In India, EPR under the Plastic Waste Management Rules drives the same move.

\textit{More layers, thinner film.} 7–11 layer dies let converters place mLLDPE, HDPE, LDPE and barrier exactly where needed and run 30–50 µm sealants.

\textit{Line automation.} Gravimetric blending, automatic die or air-ring profile control, IBC and inline dyne/gel inspection are now standard on new lamination film lines.

\textit{Lamination process shift.} Solventless polyurethane adhesives dominate; they are more sensitive to slip migration and dyne loss than solvent-based systems, which is pushing non-migratory slip and lower-slip sealant designs.

\section{Troubleshooting Lamination Defects}

Table~\ref{tab:trouble} maps common laminate defects to their film-side causes and remedies at the blown film line.

\begin{table}[!t]
\caption{Laminate Defects Traced to the PE Film}
\label{tab:trouble}
\footnotesize
\begin{tabularx}{\columnwidth}{@{}p{16mm}LL@{}}
\toprule
\textbf{Defect} & \textbf{Likely film-side cause} & \textbf{Remedy at the line}\\
\midrule
Low or falling bond & Dyne below 40; slip or antistat at treat side & Raise treatment; slip only in seal skin; non-migratory slip; laminate within 2–4 weeks\\
COF rises after cure & Erucamide migrating into adhesive and PET side & Higher initial slip, antiblock synergy, or non-migratory slip\\
Tunnelling or curl & Low BUR, gauge bands, asymmetric co-ex & BUR 2.2–3.0, oscillating haul-off, auto-profile, balanced layers\\
Gel bumps & Unmelt, oxidised gels, filler agglomerates & Finer screens, correct melt temperature, antioxidant, dry MB\\
Wrinkles, baggy edges & Gauge variation, hot nip, unstable bubble & Steady frost line, chilled air, more LDPE\\
Roll blocking & Hot winding, low antiblock, plastomer skin & Wind below 35–40 °C; raise antiblock\\
Melt fracture & mLLDPE at high shear, narrow die gap & PPA conditioning, wider gap, hotter die lip\\
\bottomrule
\end{tabularx}
\end{table}

\section{Conclusion}

LDPE remains essential in lamination film, now mainly as the 10–30\% melt-strength component that keeps a high-BUR bubble stable. A BUR of 2.0–3.0, a frost line at 2–4 die diameters and a 190–215 °C LDPE melt give balanced, flat, low-gel film suited to lamination. Metallocene PE and plastomers have cut seal initiation by 20–30 °C and allowed sealant gauge to roughly halve since the 1980s. Additives must be placed by layer, with slip, antistat and fillers kept away from the treated bonding surface. The next decade is about all-PE recyclable laminates, recycled content in core layers and PFAS-free processing, all of which put a premium on precise layer design and process control. Process windows, temperature profiles and property ranges given here are typical values for guidance and should be confirmed against resin supplier data sheets and line trials.

\section*{Acknowledgment}

The author acknowledges the use of Claude, a large language model by Anthropic, for assistance in drafting, structuring and formatting this paper. All technical content was reviewed by the author.

\begin{thebibliography}{9}
\footnotesize
\setlength\itemsep{1pt}

\bibitem{fawcett} E. W. Fawcett, R. O. Gibson, M. W. Perrin, J. G. Paton and E. G. Williams, ``Improvements in or relating to the polymerisation of ethylene,'' British Patent 471\,590, 1937.

\bibitem{cantor} K. Cantor, \textit{Blown Film Extrusion}, 2nd ed. Munich, Germany: Hanser, 2011.

\bibitem{ceflex} CEFLEX, ``Designing for a Circular Economy (D4ACE) Guidelines: Executive Summary,'' Sep. 2025. [Online]. Available: \url{https://guidelines.ceflex.eu/assets/public_docs/D4ACE_Guidelines_Executive%20Summary_Sep25.pdf}

\bibitem{pe2023} ``Flexible and recyclable monomaterial packaging meets sustainability needs,'' \textit{Plastics Engineering}, May 2023. [Online]. Available: \url{https://www.plasticsengineering.org/2023/05/flexible-and-recyclable-monomaterial-packaging-meets-sustainability-needs-000295/}

\bibitem{wh} Windmöller \& Hölscher and ExxonMobil, ``De-inkable, recyclable mono-material PE pouch with barrier properties,'' product brief. [Online]. Available: \url{https://www.wh.group/media/sustainability/23_gbd_endproduct_deinkablepouch_collab_exxon.pdf}

\bibitem{clariant} ``Clariant launches PFAS-free processing aid for polyolefin films,'' \textit{Packaging Digest}, 2025. [Online]. Available: \url{https://www.packagingdigest.com/flexible-packaging/clariant-launches-pfas-free-processing-aid-for-polyolefin-films}

\bibitem{ampacet} Ampacet Corp., ``PFAS-Free Polymer Process Aid,'' product information. [Online]. Available: \url{https://www.ampacet.com/masterbatch-products/ampacet-process/pfas-free-polymer-process-aid}

\end{thebibliography}

\end{document}
